\section{VPN wireguard}\label{vpn-wireguard}

\textbf{fonte}: \url{https://www.wireguard.com/papers/wireguard.pdf}
\textbf{vpn}: \url{https://ieeexplore.ieee.org/document/4756450}

\textbf{VPN}: e\textquotesingle{} un modo per comunicare in modo sicuro
tra membri di un gruppo attraverso l\textquotesingle uso di una
infrastruttura di comunicazione pubblica, grazie all\textquotesingle uso
di un protocollo di tunneling.

\textbf{Wireguard} e\textquotesingle{} un VPN che utilizza il concetto
di \textbf{cryptokey routing}: viene mantenuta in memoria una tabella
dove ad ogni chiave pubblica di un peer, viene associata una lista di
indirizzi ip identificati come allowed. Fornisce
un\textquotesingle interfaccia \textbf{stateless} per utilizzare il
\textbf{tunnel VPN}, ossia un\textquotesingle interfaccia di rete
virtuale, che potrebbe essere per esempio \passthrough{\lstinline!wg0!}.

\begin{quote}
per stateless, si intende il fatto che l\textquotesingle amministratore
non si deve occupare di gestire lo scambio delle chiavi, connessioni,
disconnessioni, poi che viene tutto gestito dietro le quinte.
\end{quote}

Rispetto ad OpenVPN, Wireguard e\textquotesingle{} percepito come più
sicuro ed efficiente:

\begin{itemize}
\tightlist
\item
  OpenVPN utilizza TLS ed espone molte delle sue funzionalità,
  aumentando la superficie di attacco sulla VPN.
\item
  OpenVPN utilizza interfacce TUN/TAP: agiscono in user space ma al
  costo di penalizzare di molto le performance, infatti i pacchetti
  devono essere copiati piu volte nel passagio da user space a kernel
  space.
\end{itemize}

\textbf{Scambio di chiavi}: Wireguard non utilizza un meccanismo
specifico per scambiare le chiavi. Infatti Wireguard non include un
protocollo interno per la distribuzione delle chiavi, si delega lo
scambio a meccanismi esterni e preesistenti (come SSH)

\textbf{Cosa vuol dire in pratica}? Lo scambio delle chiavi si fa per
esempio manualmente, si copia e incollano manualmente le chiavi
pubbliche nel file di configurazione di wireguard oppure si usano tool
automatizzati.

vediamo un esempio di \emph{cryptokey routing table}: !{[}{]}(Pasted
image 20261001214515.png) questa decide in pratica quali coppie
\((\text{Peer Public Key, Allowed Source IP})\) possono instradare
nell\textquotesingle interfaccia wireguard.

opzionalmente e\textquotesingle{} possibile specificare verso quale
endpoint inoltrare il pacchetto, ma almeno un peer, ad esempio il
client, deve conoscere l\textquotesingle endpoint a cui inviare i
datagrammi UDP crittografati.

\begin{quote}
l\textquotesingle{}\textbf{endpoint} identifica
l\textquotesingle indirizzo ip e la porta UDP su cui il demone WireGuard
e\textquotesingle{} in ascolto.
\end{quote}

Brevemente il flusso dei dati e\textquotesingle{} il seguente:
l\textquotesingle interfaccia \passthrough{\lstinline!wg0!} riceve un
pacchetto

\begin{itemize}
\tightlist
\item
  il client incapsula il pacchetto all\textquotesingle interno di una
  datagramma UPD il cui payload viene crittografato utilizzando
  \passthrough{\lstinline!ChaCha20Poly1305!}
\item
  il client invia il datagramma all\textquotesingle endpoint
  specificato, oppure imparato dal Wireguard in precedenza.
\item
  il server riceve il dataramma, verifica la correttezza nella
  \emph{cryptokey routing table}
\item
  il server autentica e decripta il messaggio ed estrae il pacchetto
  originario
\item
  il pacchetto originario viene inserito in una coda
  nell\textquotesingle interfaccia \(\text{wg}0\)
\end{itemize}

\section{SDN}\label{sdn}

\url{https://networklessons.com/network-automation/introduction-to-sdn-software-defined-networking}
\url{https://ridhanegara.staff.telkomuniversity.ac.id/files/2017/04/Thomas-D.-Nadeau-Ken-Gray-SDN-Software-Defined-Networks-O_039_Reilly-Media-2013.pdf}

una SDN opera su tre differenti piani:

\begin{itemize}
\tightlist
\item
  piano di controllo
\item
  paino dei dati
\item
  piano di management
\end{itemize}

Il \textbf{piano di controllo} e\textquotesingle{} responsabile per lo
scambio di informazioni sul routing, le tabelle ARP, il routing BGP,
OSPF, ed eseguire l\textquotesingle arlgoritmo STP per creare una
topologia senza cicli.

Il \textbf{piano dei dati} e\textquotesingle{} responsabile di
instradare il traffico, ossia di effettuare il forward dei pacchi, e lo
fa applicando le informazioni in possesso dal piano di controllo.
Dunque, tra le verie operazioni puo: incapsulare e decapsulre i
pacchetti, fare il match guardando MAC o indirizzi o IP, applicare le
regole NAT e oppure fare il drop dei pacchetti.

Il piano di management viene usato per accedere e gestire i device di
rete: si accede al device usando telnet, SSH, o altri servizi.
!{[}{]}(Pasted image 20261002154936.png)

Nella SDN utilizziamo un controller centrale per applicare il piano di
controllo. Il controller potrebbe essere su un host fisico oppure su una
macchina virtuale.

Cio\textquotesingle{} implica che il contorller decide come saranno
fatte le tabelle di inoltre, e gli switch ed i router si limitano ad
eseguirle. Si toglie di mezzo la complessità dovuta al setup di migliaia
di interfacce, dato che ora la gestione e\textquotesingle{}
centralizzata. Tuttavia il controller SDN e\textquotesingle{} incaricato
di fornire i dati necessari all\textquotesingle inoltro a tutti i
device. !{[}{]}(Pasted image 20261002155254.png)

Per implementare questo meccanismo ci sono due interfacce:

\begin{itemize}
\tightlist
\item
  northbound
\item
  southbound !{[}{]}(Pasted image 20261002155542.png)
\end{itemize}

\textbf{southbound}: Il controller SDN deve poter comunicare con i
device di rete in moda da programmare il piano dei dati. La southbound
interface permette la comunicazione. E\textquotesingle{}
un\textquotesingle interfaccia software che puo\textquotesingle{} essere
implementata in vari modi.

\textbf{northbound}: e\textquotesingle{} l\textquotesingle interfaccia
usata per accedere al controller SDN. Permette di applicare modifiche e
di accedere ad informazioni sul sistema

\begin{itemize}
\tightlist
\item
  Ottenere una lista di tutti i device di rete preenti
\item
  Aggiungere una nuova VLAN agli switch
\item
  Mostrare la topologia della rete
\item
  Configurare gli indirizzi ip !{[}{]}(Pasted image 20261002155846.png)
\end{itemize}

\begin{quote}
abbiamo accenato prima ad Open vSwitch. Riprendendo il discorso, il
controler SDN non fa altro che comandare gli switch gestiti da Open
vSwitch.
\end{quote}
